\documentclass[10pt]{article}
\input{../../notes-style.tex}

\usepackage{multicol}
\usepackage{array}

\title{Ch.~6 Notes\\\large Trees and Binary Search Trees}
\author{}
\date{}

\begin{document}
\maketitle
\vspace{-2em}

\begin{multicols}{2}

\subsection*{Vocabulary}
\begin{itemize}
    \item \textbf{Root} -- no parent.
    \item \textbf{Leaf} -- no children.
    \item \textbf{Depth} of a node: edges from root. Root depth = 0.
    \item \textbf{Height} of a tree: max depth. Empty tree $= -1$.
    \item \textbf{Level} $k$: all nodes at depth $k$.
    \item \textbf{Full}: every node has 0 or 2 children.
    \item \textbf{Complete}: all levels full except possibly the last,
    which fills left-to-right.
    \item \textbf{Perfect}: all levels full.
\end{itemize}
Perfect binary tree of height $h$: $2^{h+1} - 1$ nodes; $2^h$ leaves.

\subsection*{Node type}
\begin{lstlisting}[language=C++]
struct TreeNode {
    int key;
    TreeNode* left  = nullptr;
    TreeNode* right = nullptr;
    // optional: TreeNode* parent;
};
\end{lstlisting}

\subsection*{BST invariant}
For every node $n$:
\begin{itemize}
    \item every key in the left subtree is $< n.key$;
    \item every key in the right subtree is $> n.key$.
\end{itemize}
(Duplicates disallowed by convention; if allowed, pick one side and
document it.)

\subsection*{BST search}
\begin{lstlisting}[language=C++]
TreeNode* bstSearch(TreeNode* root, int k) {
    while (root && root->key != k)
        root = (k < root->key) ? root->left
                               : root->right;
    return root;
}
\end{lstlisting}
Cost: $O(h)$ -- $O(\log n)$ balanced, $O(n)$ degenerate.

\subsection*{BST insert}
Walk as if searching; when you fall off a null pointer, that's where
the new node goes.
\begin{lstlisting}[language=C++]
void bstInsert(TreeNode*& root, int k) {
    if (!root) { root = new TreeNode{k}; return; }
    if (k < root->key)  bstInsert(root->left,  k);
    else if (k > root->key) bstInsert(root->right, k);
    // else duplicate: policy decision
}
\end{lstlisting}

\columnbreak

\subsection*{BST remove: four cases}
\begin{enumerate}
    \item \textbf{Leaf.} Delete the node; parent's pointer $\to$ null.
    \item \textbf{Only left child.} Parent's pointer $\to$ that child.
    \item \textbf{Only right child.} Parent's pointer $\to$ that child.
    \item \textbf{Two children.} Find \emph{in-order successor}
    (leftmost of right subtree). Copy its key into the node to be
    deleted. Remove the successor (it has at most one child $\to$ case
    1 or 3).
\end{enumerate}
Why successor? It's guaranteed $>$ everything in left subtree and $<$
everything in right subtree -- preserves the BST invariant.

\subsection*{Traversals}
\begin{tabular}{@{}l|l|l@{}}
\textbf{Order}         & \textbf{Visit pattern}    & \textbf{Use}\\
\hline
Pre-order              & root, L, R               & copy/serialize\\
In-order               & L, root, R               & sorted output\\
Post-order             & L, R, root               & delete tree
\end{tabular}
All three are $O(n)$; recursion depth $O(h)$.

\subsection*{In-order, iterative (useful!)}
\begin{lstlisting}[language=C++]
void inorder(TreeNode* root) {
    std::stack<TreeNode*> s;
    TreeNode* cur = root;
    while (cur || !s.empty()) {
        while (cur) {
            s.push(cur); cur = cur->left;
        }
        cur = s.top(); s.pop();
        visit(cur);
        cur = cur->right;
    }
}
\end{lstlisting}

\subsection*{Set-ADT order ops on a BST ($O(h)$ each)}
\texttt{find\_min} -- walk left until null.\\
\texttt{find\_max} -- walk right until null.\\
\texttt{find\_next(x)} (in-order successor):
\begin{itemize}
    \item if \texttt{x->right} non-null: leftmost of right subtree;
    \item else: walk up via parent until you come from the left.
\end{itemize}
\texttt{find\_prev} -- symmetric. \emph{This is what hash tables
can't do} -- they cost $O(n)$ for any of these.

\subsection*{Height vs.\ balance}
Best: $h = \lfloor \log_2 n \rfloor \to O(\log n)$ ops.\\
Worst: ascending/descending insert sequence $\to$ ``rope'' tree,
$h = n - 1$, ops become $O(n)$.\\
\textbf{Random insertion order:} expected $h = O(\log n)$ (CLRS
Thm~12.4) -- a vanilla BST built from random data already behaves
like a balanced one on average. \emph{Production uses red-black
(ch.~9) anyway for the worst-case guarantee.}
\textbf{Balance = the difference.}
Balance factor of a node = $\text{height}(\text{left}) -
\text{height}(\text{right})$. AVL enforces $|bf| \leq 1$.

\subsection*{BST vs.\ hash table vs.\ sorted array}
\begin{tabular}{@{}l|lll@{}}
                  & \textbf{BST} & \textbf{hash} & \textbf{sorted arr}\\
\hline
find              & $O(\log n)$  & $O(1)^*$      & $O(\log n)$\\
insert            & $O(\log n)$  & $O(1)^*$      & $O(n)$\\
sorted iter       & $O(n)$       & $O(n \log n)$ & $O(n)$\\
range query       & yes          & no            & yes\\
min/max           & $O(\log n)$  & $O(n)$        & $O(1)$
\end{tabular}
$^*$ expected. BST numbers assume balance.

\subsection*{Parent pointers: why/why not}
\textbf{With parent:} makes remove cleaner, supports iterator-style
traversal, required for AVL / red-black rotations. Costs one extra
pointer per node.\\
\textbf{Without parent:} smaller nodes; reach prev/next via recursion
stack or reference-to-pointer idiom.

\subsection*{Top gotchas}
\begin{itemize}
    \item Quoting BST ops as $O(\log n)$ without saying \emph{balanced}.
    \item Removing with two children by replacing with the \emph{root}
    of a subtree (invalid) instead of the in-order successor.
    \item Forgetting that insertion order determines tree shape;
    sorted input $\to$ linear chain.
    \item Using \texttt{TreeNode*} (value) instead of
    \texttt{TreeNode*\&} (reference) when a function needs to alter
    the parent's pointer.
    \item Memory leaks from not deleting the tree (use post-order).
\end{itemize}

\end{multicols}
\end{document}
